Man vergleicht die Strecken $\Delta s$, die in aufeinanderfolgenden gleich langen Zeitintervallen $\tau = \tZbO$ zurückgelegt werden; $\Delta s/\tau$ ist die mittlere Geschwindigkeit im jeweiligen Intervall. Bei einer gleichmässig beschleunigten Bewegung nimmt die Geschwindigkeit in jedem Intervall um gleich viel zu, die Strecken $\Delta s$ wachsen also jedes Mal um denselben Betrag $a\,\tau^2$.

\begin{center}
\begin{tabular}{l|ccccc}
Intervall & 1 & 2 & 3 & 4 & 5 \\
\hline
$\Delta s_1$ & \dAaP{0}{2} & \dAbP{0}{2} & \dAcP{0}{2} & \dAdP{0}{2} & \dAeP{0}{2} \\
Zunahme & \eAaP{0}{2} & \eAbP{0}{2} & \eAcP{0}{2} & \eAdP{0}{2} &  \\
\hline
$\Delta s_2$ & \dBaP{0}{2} & \dBbP{0}{2} & \dBcP{0}{2} & \dBdP{0}{2} & \dBeP{0}{2} \\
Zunahme & \eBaP{0}{2} & \eBbP{0}{2} & \eBcP{0}{2} & \eBdP{0}{2} &  \\
\hline
$\Delta s_3$ & \dCaP{0}{2} & \dCbP{0}{2} & \dCcP{0}{2} & \dCdP{0}{2} & \dCeP{0}{2} \\
Zunahme & \eCaP{0}{2} & \eCbP{0}{2} & \eCcP{0}{2} & \eCdP{0}{2} &  \\
\end{tabular}
\end{center}

\begin{itemize}
\item Bewegung 1 ist $\result{\text{gleichmässig beschleunigt}}$: Die Strecken pro Sekunde nehmen jedes Mal um \eAaP{0}{2} zu. Die Beschleunigung ist
\begin{align}
 a &= \aAF = \frac{\dAe-\dAa}{4\cdot\left(\tZb\right)^2} = \aA \approx \aAP{0}{2}
\end{align}
\item Bewegung 2: Der Körper wird immer schneller, aber die Zunahme der Strecken pro Sekunde wird kleiner. Die Bewegung ist beschleunigt, die Beschleunigung nimmt jedoch ab (wie bei einem Auto, das sich seiner Höchstgeschwindigkeit nähert).
\item Bewegung 3: Der Körper wird immer schneller, und die Zunahme der Strecken pro Sekunde wird grösser. Die Beschleunigung nimmt mit der Zeit zu.
\end{itemize}
